% Proof of Kamp III.1, retaining all quantifier-position and empty-region cases.

\section{Kamp's first-order normal form over linear orders}
\label{sec:annex-p09-kamp-iii1}

\begin{annexitem}{annex:p9.kamp.iii1}{thm:p09-kamp-iii1}{The linear-order normal-form construction}
Use the language, prefixes, components and standard formulas of
Definitions~\ref{def:p09-kamp-iii1-prefix}--\ref{def:p09-kamp-iii1-standard}.
This is the construction of Kamp's Theorem~III.1, with explicit
capture-avoiding substitutions, constants and witness-existence conditions
\cite[ch.~III, pp.~43--56]{Kamp1968}.

\begin{proof}
We first construct the starred normal form and prove its additional
predicate-support property.  Until the paragraph converting it back to the
source grammar, ``standard formula'', ``component'', and
\(\mathcal C_\pi\) refer to their starred versions, and the component
names \(\operatorname{At},\operatorname{Bef},\operatorname{Aft},
\operatorname{Bet}\) suppress that star.  Prefixes and their equality
abbreviation are unchanged.  Thus the occurrences of \(\top_x,\bot_x\)
in this induction are explicitly in \(\operatorname{At}^{*}(x)\), not
in Kamp's literal \(\operatorname{At}(x)\).

\paragraph{Order types and finite Boolean rearrangements.}
Under a linear-order assignment, coincidence partitions a finite nonempty
variable set \(X\), and strict order uniquely orders the resulting blocks.
The standard prefix \(\pi\) for that ordered partition is true: its chain
records the adjacent blocks and its equalities identify all other variables
with their representatives.  Conversely those conjuncts, transitivity and
linearity force precisely that ordered partition.  Therefore the finitely
many standard prefixes on \(X\) are exhaustive and pairwise incompatible.
Each is consistent, since it is realized by the corresponding finite chain.
It follows also that any consistent temporal prefix is equivalent to the
disjunction of exactly those standard prefixes whose order types satisfy
it.  This proves the prefix facts used below rather than assuming that a
prefix describes a dense or endpointless order.

The only Boolean rearrangements needed are finite distribution, regrouping
of conjunctions and disjunctions, and deletion of duplicate conditions.
For a fixed prefix \(\pi\), conjunctions and disjunctions of standard
clauses with that prefix can be distributed back into standard clauses with
the same prefix: every resulting local conjunct remains in
\(\mathcal C_\pi\).  Clauses from two different complete prefixes have
false conjunction.  Missing local conjunctions can always be filled by
\(\top\) at the first representative, and a false branch by \(\bot\)
there.  Thus these finite rearrangements preserve standard form and the
named free variables.

\paragraph{Quantifier-free basis.}
First rename bound variables apart and put the input in prenex form, using
the usual Boolean and quantifier equivalences.  These transformations
preserve the free variables; moving a quantifier across a formula in which
its variable is not free is legitimate, and vacuous quantifiers can be
removed because the order is nonempty.  Induct on the number of quantifiers
in this prenex input, retaining the entire set of free variables at each
step.

For the basis, put the quantifier-free formula in disjunctive normal form.
In each conjunction, separate unary literals from order literals.  Over
\(\mathsf{Lin}\) the relevant rewrites are
\[
 \neg(u\doteq v)\longleftrightarrow(u<v\lor v<u),\qquad
 \neg(u<v)\longleftrightarrow(u\doteq v\lor v<u).
\]
Distribute again, or equivalently enumerate the finite ordered partitions
of the free-variable set and retain just those satisfying the order part.
For every retained prefix, replace the variable in each unary literal by
the representative of its equality block.  Coincidence of the assigned
points justifies this replacement for every unary interpretation.
Conjoin literals at each representative, producing
\(\operatorname{At}\) components.  A contradictory unary conjunction
remains an allowed false component; if the order part is inconsistent no
order type is retained.  If the order part or the unary part is empty, use
respectively all compatible prefixes or \(\top\) at the first
representative.  If no clause survives, use the explicit always-false
standard formula described in Definition~\ref{def:p09-kamp-iii1-standard}.
This covers also formulas containing no unary predicate.

\paragraph{Existential step: the two equality cases.}
Suppose that the normalized inner formula has the form
\[
 N(x,X)=\bigvee_i\left(\pi_i\land\bigwedge_j\theta_{ij}\right),
 \qquad X\ne\varnothing,
\]
where the \(\pi_i\) are standard on \(X\cup\{x\}\).
Since existential quantification distributes over finite disjunction, it
suffices to normalize
\(\exists x\,(\pi\land\bigwedge_j\theta_j)\).
If the original inner formula did not have \(x\) free, first remove the
vacuous quantifier and use its normal form on \(X\).

There are two distinct equality configurations in the source convention.
If \(x\) is not essential, its block has representative \(y\), and the
only prefix conjunct containing \(x\) is \(y\doteq x\).
No local component has \(x\) free, since only essential variables name
component endpoints.  Delete \(y\doteq x\) and use the standard prefix
\(\rho\) on \(X\).  When only \(y\) remains, this is
\(y\doteq y\).  The result is
\[
 \rho\land\bigwedge_j\theta_j.
\]
Indeed the forward implication follows by deleting the equality, and for
the reverse implication assign \(x\) the already existing point assigned
to \(y\).  No separate existence hypothesis is required in an equality
case.

If \(x\) is essential but shares its block with other variables, let
\(z\) be the first remaining variable of that block.  Replace free
occurrences of \(x\) by \(z\) in the prefix and all its components,
delete the redundant \(z\doteq z\) unless \(X=\{z\}\), and write
the resulting prefix canonically as \(\rho\).  The new components
\(\theta_j[z/x]\) have exactly the permitted endpoints for \(\rho\).
To see this for nested components, rename all their bound variables fresh
first, then inspect the three recursive constructors: replacing an endpoint
changes the same endpoint in its guard and its constituent components and
leaves the constructor in the corresponding class.  Thus the result is
\[
 \rho\land\bigwedge_j\theta_j[z/x].
\]
Equivalence again follows by the equality \(x\doteq z\), and the reverse
direction chooses its forced witness.  These are respectively cases 2(a)
and 2(b) of the source proof.

\paragraph{Existential step: the three strict positions.}
The remaining case has \(\{x\}\) as an entire equality block.  Remove
that block and call the resulting standard prefix on \(X\) \(\rho\).
The position condition \(\beta\), the parts of the prefix removed or
joined, and the possible local components containing \(x\) are exactly
as follows:
\[
\begin{array}{c|c|l}
 \text{position of }x&\beta&\text{components with }x\text{ free}\\ \hline
 \text{before the first }y&x<y&
 \operatorname{At}(x),\ \operatorname{Bef}(x),\ \operatorname{Bet}(x,y)\\
 \text{after the last }y&y<x&
 \operatorname{At}(x),\ \operatorname{Aft}(x),\ \operatorname{Bet}(y,x)\\
 \text{between adjacent }y,z&y<x\land x<z&
 \operatorname{At}(x),\ \operatorname{Bet}(y,x),\ \operatorname{Bet}(x,z)
\end{array}
\]
In the first two rows remove the one strict comparison containing \(x\).
In the third replace \(y<x\land x<z\) by \(y<z\).
If one variable remains, insert its self-coincidence prefix.  In every row
\(\pi\) is equivalent to \(\rho\land\beta\).

Write \(I\) for the conjunction of the local components with no free
\(x\), and \(D\) for the conjunction of the others.  Every component of
\(I\) is allowed for \(\rho\): deleting the singleton block affects
only the two neighboring cells, or the one exterior cell and its neighbor.
Every component of \(D\) lies in the indicated row of the display.
The exact transformation is
\[
 \exists x\,(\pi\land I\land D)
 \longleftrightarrow
 \rho\land I\land\exists x\,(\beta\land D).
\]
If there is no component in \(D\), put \(D=\top_x\); the last
existential is retained.  The resulting quantified component belongs to
\(\operatorname{Bef}(y)\), \(\operatorname{Aft}(y)\), or
\(\operatorname{Bet}(y,z)\), respectively.  More explicitly, choose a
finite stage \(n\) containing all its finitely many non-\(\operatorname{At}\)
constituents; the new component is in stage \(n+1\) by Definition~
\ref{def:p09-kamp-iii1-components}.

The equivalence holds because \(\rho\) and \(I\) are independent of
the quantified assignment, and both sides ask for the same point satisfying
\(\beta\land D\).  If the left ray, right ray, or intervening gap is
empty, both sides are false.  Thus the first and last points of a finite
order and adjacent points of a discrete order are included.  These are all
three subcases of source case 2(c), with their existence conditions made
explicit.  Disjoining the results proves the existential induction step.

\paragraph{Universal step: restrict prefixes and retain each position.}
For \(\forall x\,N(x,X)\), with \(X\ne\varnothing\), first handle a
vacuous quantifier as above.  Each inner standard clause can be put in the
form
\[
 \rho\land\beta\land\bigwedge_j\theta_j,
\]
where \(\rho\) is standard on \(X\) and \(\beta\) records just the
position of \(x\).  The source's four prefix operations are:

\begin{enumerate}
\item If \(x\) is an essential member of a nonsingleton block, choose its
      first remaining member \(y\), replace \(x\) by \(y\) in all local
      components and in the restricted prefix, and use
      \(\beta=(x\doteq y)\).
\item If \(x\) is not essential, remove its one equality
      \(y\doteq x\), retain the already \(x\)-independent components,
      and use \(\beta=(y\doteq x)\).
\item If \(x\) is a singleton first or last block, delete its one strict
      comparison and take it as \(\beta\).
\item If \(x\) is a singleton interior block, replace its two comparisons
      \(y<x\land x<z\) by \(y<z\), taking the two original comparisons
      together as \(\beta\).
\end{enumerate}
The one-variable residual prefix is again a self-coincidence.  In the
equality cases the substituted local conditions no longer contain \(x\)
free.  These operations are justified by the same order-type and substitution
arguments as in the existential step; unlike that step, universal
quantification cannot simply be distributed over the original disjunction.

Fix now a residual prefix \(\rho\) with essential variables
\(y_1,\ldots,y_s\).  All positions of \(x\) are
\[
\begin{gathered}
 x<y_1,\qquad y_k<x\land x<y_{k+1}\ (1\leq k<s),
 \qquad y_s<x,\\
 x\doteq y_k\ (1\leq k\leq s).
\end{gathered}
\]
In the last row reverse the displayed equality when the surviving
representative precedes \(x\) in variable-name order.  This recovers the
two equality orientations of the source; the meaning is the same.
There are \(2s+1\) positions.  Given \(\rho\), they are exhaustive
and mutually exclusive by linearity.  An equality position always has its
named point; a strict position may have no point at all.

Group all original clauses first by their residual \(\rho\), then by
their position \(\beta\).  Let
\[
 H_{\rho,\beta}:=\bigvee_{i\text{ in this group}}
                       \bigwedge_j\theta_{ij}.
\]
For an absent strict-position group, use \(H_{\rho,\beta}=\bot_x\).
For an absent equality group \(x\doteq y_k\), use
\(H_{\rho,\beta}=\bot_{y_k}\).
Include every standard prefix \(\rho\) on \(X\), even if it had no
original group.  Exact partitioning then gives
\[
 \forall x\,N(x,X)\longleftrightarrow
 \bigvee_\rho\left(
   \rho\land\bigwedge_{\beta}
          \forall x\,(\beta\longrightarrow H_{\rho,\beta})
 \right).
\]
For any fixed assignment to \(X\), exactly one \(\rho\) is true;
for each assignment to \(x\), exactly one of that prefix's positions is
true.  The right-hand side therefore tests precisely the same matrix at
every point as the left-hand side.  In particular a missing strict-position
group requires that region to be empty; a missing equality group makes the
whole branch false because its named point exists.  These are the missing
\(\gamma_f\) positions in source equation (10), not deletable disjuncts.

\paragraph{Universal step: conjunctive normal form inside a position.}
Use finite Boolean distribution, treating the components as propositional
atoms, to write each nonconstant \(H_{\rho,\beta}\) as
\(\bigwedge_m\bigvee_p\eta_{mp}\).  Every \(\eta_{mp}\) is one
of the group's original local components; no negation of a component is
introduced.  Constants can be kept as their one-component clauses.
Distribute the universal quantifier over the finite conjunction.  It remains
to normalize
\[
 \rho\land\forall x\left(\beta\longrightarrow
                              \bigvee_p\eta_p\right).
\]

If \(\beta\) is an equality position, all the \(\eta_p\) have no
free \(x\), after the prior representative substitution.  The point
\(y_k\) realizes this position, so the formula is equivalent to
\[
 \bigvee_p(\rho\land\eta_p).
\]
All these are standard clauses for \(\rho\).  A missing equality group
has the false component \(\bot_{y_k}\) and gives a false branch by the
same argument.  This deals with both possible orientations of equality.

Suppose instead that \(\beta\) is strict.  Separate the disjunction into
\(A\lor B\), where \(A\) contains the components with no free \(x\)
and \(B\) those with free \(x\).  The precise three cases are:

\begin{enumerate}
\item If all components contain \(x\), the formula is already one standard
      clause: its universal component has the appropriate
      \(\operatorname{Bef}\), \(\operatorname{Aft}\), or
      \(\operatorname{Bet}\) form.
\item If none contains \(x\), replace the formula by
      \[
       \bigvee_p(\rho\land\eta_p)
       \ \lor\
       \left(\rho\land\forall x\,(\beta\longrightarrow\bot_x)\right).
      \]
      If a component \(\eta_p\) is true, every point in the region satisfies
      the consequent.  If none is true, the universal condition holds
      exactly when the region is empty.  The last disjunct records this
      possibility and is an allowed universal component.
\item If both kinds occur, replace the formula by
      \[
       \bigvee_{\eta_p\text{ in }A}(\rho\land\eta_p)
       \ \lor\
       \left(\rho\land\forall x\left(
          \beta\longrightarrow\bigvee_{\eta_p\text{ in }B}\eta_p
       \right)\right).
      \]
      The disjunction \(A\) has a fixed truth value while \(x\) varies.
      If it is true the condition holds everywhere; if false the original
      condition reduces to the remaining universal.  This proves the
      equivalence also on an empty region.
\end{enumerate}

For \(x<y_1\), the \(x\)-dependent constituents are in
\(\operatorname{At}(x)\cup\operatorname{Bef}(x)
\cup\operatorname{Bet}(x,y_1)\), so each newly formed universal
component is in \(\operatorname{Bef}(y_1)\).
For \(y_s<x\), they are in
\(\operatorname{At}(x)\cup\operatorname{Aft}(x)
\cup\operatorname{Bet}(y_s,x)\), giving
\(\operatorname{Aft}(y_s)\).
For \(y_k<x<y_{k+1}\), they are in
\(\operatorname{At}(x)\cup\operatorname{Bet}(y_k,x)
\cup\operatorname{Bet}(x,y_{k+1})\), giving
\(\operatorname{Bet}(y_k,y_{k+1})\).
The falsehood \(\bot_x\) in an absent group is itself an
\(\operatorname{At}(x)\) component, so exactly the same constructor
handles every missing region.  In each case take the maximum of the finitely
many stages of the constituents, then the next stage, to verify membership.
The \(x\)-independent constituents are allowed local components for
\(\rho\), as established in the prefix restriction step.

Finally conjoin the finitely many resulting formulas for all positions and
all their conjunctive-normal-form clauses, distributing back to disjunctive
form.  Every clause retains the same \(\rho\) and only allowed local
components.  Disjoin over \(\rho\).  The initial finite Boolean
observation proves that the result is a standard formula on \(X\).
This completes every universal case and hence the quantifier induction.

\paragraph{Effectivity, free variables and support.}
There are finitely many ordered partitions of each finite variable set.
The two Boolean conversions, partition enumeration, choice of representatives,
fresh-variable renaming and finite grouping operations are syntactic finite
procedures.  An induction step normalizes one outer quantifier of an already
normalized inner formula; it does not recursively renormalize the new
regional components.  Thus the original finite quantifier prefix is
exhausted.  The recursive component definitions certify the newly constructed
formulas at finite stages, regardless of the size of the order.
Every output clause retains a prefix on all the required free variables.
Only existing unary literals are moved or repeated; the constants and
coincidence abbreviation use only \(<\).  Neither new free variables nor
new unary predicates are introduced.  This proves the effective starred
normal-form assertion and the additional handbook support property, without
asserting an output-size or complexity bound.

\paragraph{Conversion to the literal source grammar.}
Restore the stars on this output, calling it \(N^*_{\varphi}\), and
apply \(\varepsilon\).  In an \(\operatorname{At}^{*}(x)\) slot,
a unary-literal conjunction is already a source \(\operatorname{At}(x)\)
formula, while \(\top_x\) and \(\bot_x\) become respectively the
allowed \(T^K_x\) and \(F^K_x\).  Their truth values agree on every
linear order, for every interpretation of \(P_0\).
Induct on the stage of a regional component.  The conversion preserves its
guard and outer quantifier; the induction hypothesis places every converted
constituent in the corresponding source class at the previous stage and
preserves its truth for every assignment.  Each converted regional component
is consequently a literal source component with the same truth condition.
The finite clause structure and all prefix variables are unchanged, so
\(\varepsilon(N^*_{\varphi})\) is a source-standard formula on exactly
the original free variables and is equivalent to \(\varphi\).
Only the inessential \(P_0\) encodings of constants can enlarge unary
support in this final conversion.  This proves the source Theorem~III.1
separately from the starred support property.  In the remaining paragraph,
``standard'' again means the unstarred source notion.

\paragraph{Sentence corollary.}
Proceed by structural decomposition of a sentence.  The immediate arguments
of a Boolean connective in a sentence are themselves sentences, so keep
that Boolean connective and apply the decomposition to them.  At a closed
quantified formula \(Qx\,\psi\), either \(x\) is free in \(\psi\),
in which case it is its only free variable, or \(\psi\) is already a
sentence.  In the first case the proved theorem gives a standard
\(N(x)\) equivalent to \(\psi\), and \(Qx\,N(x)\) is a standard
sentence for either \(Q=\exists\) or \(Q=\forall\).  In the second
case remove the vacuous quantifier, valid on a nonempty order, and continue
with the smaller sentence \(\psi\).  The language has no closed atomic
formulas, since its atomic formulas contain variables and it has no constant
symbols.  If Boolean truth or falsehood abbreviations have been introduced,
represent them by \(\exists x\,((x\doteq x)\land T^K_x)\) and
\(\forall x\,((x\doteq x)\land F^K_x)\), respectively.
This finite decomposition gives a Boolean combination of standard sentences,
exactly the scope of Kamp's Corollary~1 on printed p.~56.
\end{proof}
\end{annexitem}
